\documentclass[11pt,a4paper]{amsart}

\usepackage[utf8]{inputenc}
\usepackage[T1]{fontenc}
\usepackage[ngerman]{babel}
\usepackage{amsmath,amssymb,amsthm}
\usepackage{array,booktabs,longtable}
\usepackage{hyperref}

\newtheorem{theorem}{Satz}[section]
\newtheorem{proposition}[theorem]{Proposition}
\newtheorem{lemma}[theorem]{Lemma}
\newtheorem{corollary}[theorem]{Korollar}
\newtheorem{definition}[theorem]{Definition}
\newtheorem{remark}[theorem]{Bemerkung}
\newtheorem{openproblem}[theorem]{Offenes Problem}

\newcommand{\RH}{\mathrm{RH}}
\newcommand{\SPC}{\mathrm{SPC}}
\newcommand{\LFF}{\mathrm{LFF}}

\title[P07 --- Weil-Form Statistics]{P07 --- Weil-Form Statistics:\\
Correlation Channels, Form-Factor Limits and Herglotz Interfaces}
\author{Objekt-X-Programm}
\date{2026-08-08}

\begin{document}

\begin{abstract}
Wir organisieren die quadratische Weilform-Statistik in fünf konzeptionellen
Schichten: den quadratischen Pivot und seinen Positivitätsbereich, die
Korrelationsstruktur mit Bochner-Lift und Skalenanalyse, den
Formfaktor-/Rampenkanal mit Goldston--Montgomery-Varianz und offenem Symboltest,
das Herglotz-Interface zwischen linearer Distribution und quadratischer
Weil-Geometrie sowie die konditionale Jacobi-/Herglotz-Realisierungsarchitektur.
Das Manuskript konsolidiert ausschließlich den auditierten Endstand von
NEU-091--120. Es ist ein internes SYN-Manuskript; eine Behauptung der
Publikationsreife ist damit nicht verbunden.
\end{abstract}

\maketitle

% ----------------------------------------------------------------
\section{Quadratischer Pivot und Positivitätsrahmen}
% ----------------------------------------------------------------

\begin{theorem}[Determinanten-No-Go]\label{thm:det-no-go}
Im festen $z$-Regime gilt
\[
  D_N(z)\xrightarrow[N\to\infty]{}e^{-\gamma^2/4}.
\]
Der Grenzwert ist damit $z$-unabhängig. Folglich entsteht in diesem Regime
kein direkter nichttrivialer $\xi$-Anschluss über den Determinantenweg.
\end{theorem}

\begin{remark}[Status]
Theorem~\ref{thm:det-no-go} ist das gesicherte Negativresultat aus NEU-091,
Status $\checkmark[M]_{\rm neg}$.
\end{remark}

\begin{definition}[Quadratischer Pivot]\label{def:quadratic-pivot}
Unabhängig vom Determinantenansatz wird
\[
  Q_N(\varphi)
  :=\gamma^2\sum_{r,n}\Lambda(n)^2\,W_N(r,n)\,\varphi(r,n)
\]
als separates quadratisches Mangoldt-Objekt eingeführt.
\end{definition}

\begin{theorem}[Testkegel-Positivität]\label{thm:test-cone}
Für $\phi$ im vorgesehenen Testkegel $\mathcal C$ gilt
\[
  Q_N[\phi]\ge 0.
\]
\end{theorem}

\begin{openproblem}[Bilinearer Lift]\label{op:bilinear-lift}
Offen bleibt, ob der echte bilineare Lift $B_N(f,g)$ mit Kreuztermen
$\Lambda(m)\Lambda(n)$ auf dem vollen Schwartzraum $\mathcal S$
positiv-semidefinit ist.
\end{openproblem}

% ----------------------------------------------------------------
\section{Korrelationsstruktur und Skalenanalyse}
% ----------------------------------------------------------------

\begin{proposition}[Abstrakte PSD-Architektur]\label{prop:psd-architecture}
Sei $\rho_N=(\rho_N(a,b))_{a,b}$ eine positiv-semidefinite Matrix mit
$\rho_N(a,a)=1$. Dann ist
\[
  K_N(a,b):=\sqrt{\kappa_a\kappa_b}\,\rho_N(a,b)
\]
positiv-semidefinit.
\end{proposition}

\begin{definition}[Logarithmischer Bochner-Kandidat]\label{def:bochner-candidate}
Als logarithmischer Kandidat wird
\[
  \rho_N((r,m),(r,n))
  :=\Psi_N\!\left(\log\frac{m}{n}\right)
\]
betrachtet, wobei $\Psi_N$ positiv-definit im Bochner-Sinn ist und
\[
  \Psi_N(t)=\int_{\mathbb R}e^{it\xi}\,d\nu_N(\xi)
\]
für ein positives Maß $\nu_N$ besitzt.
\end{definition}

\begin{openproblem}[Kanonische Wahl von $\rho_N$]\label{op:rho-canonical}
Die abstrakte PSD-Architektur bestimmt noch keine kanonische Wahl von $\rho_N$.
\end{openproblem}

\begin{theorem}[Skalentriage]\label{thm:scale-triage}
Mit $M_N$ als Kurzintervallparameter und $T=M_N^\theta$ ergeben sich drei
Regime:
\begin{center}
\small
\begin{tabular}{>{\raggedright\arraybackslash}p{0.19\textwidth}
                >{\raggedright\arraybackslash}p{0.34\textwidth}
                >{\raggedright\arraybackslash}p{0.35\textwidth}}
\toprule
Regime & Verhalten & Bedeutung\\
\midrule
$T=O(1)$
& Rang-eins-/Grobmittelungs-Kollaps
& zu grob für einen Weil-Anschluss\\
$1\ll T\ll M_N$
& nichttrivialer Kandidatenbereich
& Weil-Tauglichkeit offen\\
$T\gg M_N$
& Diagonal-Kollaps
& zu fein; kein Weil-Anschluss\\
\bottomrule
\end{tabular}
\normalsize
\end{center}
\end{theorem}

\begin{proposition}[Feinstruktur des Zwischenregimes]\label{prop:fine-structure}
Im Zwischenregime $1\ll T\ll M_N$ ist der auditierte Endstand:
\begin{enumerate}
  \item Der Singulärserien-Hauptterm steht unter der
        Hardy--Littlewood-Siebheuristik und ist daher \emph{CONDITIONAL}.
  \item Die Restdichte $\Delta_N$ und der Übergang zum Shift-Spektrum bleiben
        als gesicherter Bestandteil der Feinstrukturanalyse erhalten.
  \item Die Weil-Tauglichkeit des Zwischenregimes bleibt offen.
\end{enumerate}
\end{proposition}

% ----------------------------------------------------------------
\section{Formfaktor-/Rampenkanal}
% ----------------------------------------------------------------

\begin{definition}[Dyadische Kurzintervallvarianz]\label{def:gm-variance}
Mit $\psi(x)=\sum_{n\le x}\Lambda(n)$ sei
\[
  \mathcal V(M,H)
  :=\frac1M\int_M^{2M}
  \bigl(\psi(x+H)-\psi(x)-H\bigr)^2\,dx.
\]
\end{definition}

\begin{theorem}[Goldston--Montgomery-Varianzasymptotik, CONDITIONAL]
\label{thm:gm}
Unter RH gilt für jedes feste $\varepsilon>0$, uniform im entsprechenden
Goldston--Montgomery-Bereich
\[
  1\le H\le M^{1-\varepsilon},
\]
die konditionale Äquivalenz
\[
  \SPC
  \quad\Longleftrightarrow\quad
  \mathcal V(M,H)\sim H\log(M/H).
\]
Für $H\asymp M$ gelten abweichende Skalenregimes.
\end{theorem}

\begin{remark}[Status]
Theorem~\ref{thm:gm} ist \emph{CONDITIONAL}; die Varianzasymptotik folgt nicht
aus RH allein. Dies ist der synchronisierte Endstand von NEU-101 und P07
\cite{GoldstonMontgomery1987}.
\end{remark}

\begin{corollary}[Selbstdualer Testwert]\label{cor:self-dual}
Am kanonischen selbstdualen Punkt $H=\sqrt M$ gilt konditional unter RH und SPC
\[
  \mathcal V(M,\sqrt M)
  \sim \sqrt M\log\sqrt M
  =\frac12\sqrt M\log M.
\]
\end{corollary}

\begin{definition}[Lokales Formfaktor-Ersatzobjekt]\label{def:local-form-factor}
Das unnormalisierte lokale Objekt $\mathcal P^{\rm unf}_{N,H}$ ersetzt das
verworfene global normierte Objekt $\mathcal S_{N,H}$. Sein gesicherter Umfang
ist nur partiell, entsprechend $\checkmark[M]_{\rm part}$ in NEU-104.
\end{definition}

\begin{theorem}[Rampenasymptotik, einseitig]\label{thm:ramp-one-way}
Im auditierten Setup gilt
\[
  \LFF=1
  \quad\Longrightarrow\quad
  \text{universelle Rampenmassenasymptotik}.
\]
Die Umkehrung ist nicht bewiesen.
\end{theorem}

\begin{theorem}[LFF-Typisierungswarnung]\label{thm:lff-firewall}
Der lokale Formfaktor allein konstruiert oder identifiziert
$Q_{\rm Weil}$ nicht.
\end{theorem}

\begin{openproblem}[Lokaler Symboltest]\label{op:symbol-test}
Der Hauptsymboltest
\[
  \sigma_{\rm loc}(Q_{\rm Weil})\stackrel{?}{=}|\alpha|
\]
bleibt offen. Weder Gleichheit noch Ungleichheit ist im auditierten Block
bewiesen.
\end{openproblem}

% ----------------------------------------------------------------
\section{Linearer Herglotz-Kanal versus quadratische Weil-Geometrie}
% ----------------------------------------------------------------

\begin{remark}[Verbindliche Fourier- und Weil-Konventionen]
\label{rem:p02-conventions}
Die Fourierkonvention und die vollständige hermitesche Weil-Typisierung werden
aus P02 übernommen. Insbesondere gilt
\[
  \widehat f(t)=\int_{\mathbb R}f(u)e^{itu}\,du,
\]
und die quadratische Ebene bleibt typgetrennt als
\[
\begin{aligned}
  (a,b)&\longmapsto g_{a,b}\longmapsto h_{a,b}\\
       &\longmapsto B_W(a,b).
\end{aligned}
\]
Der adelische Port $R_{\rm PW}$ sowie die komponentenweise Definition von
$B_W$ werden daher nicht erneut dupliziert. P03 dient hier nur dem
Stil- und Konventionsabgleich.
\end{remark}

\begin{definition}[Arithmetische Herglotzfunktion in Nevanlinna-Form]
\label{def:m-arith}
Unter RH sei
\[
  \Gamma:=\{\gamma\in\mathbb R:\xi(1/2-i\gamma)=0\},
  \qquad
  \mu_{\rm arith}:=\sum_{\gamma\in\Gamma}m_\gamma\delta_\gamma.
\]
Das Maß $\mu_{\rm arith}$ ist das reine Nullstellenmaß; Gamma-, Pol- und
Primterme sind keine zusätzlichen Bestandteile dieses Herglotz-Spektralmaßes.
Die zugehörige arithmetische Funktion wird in der Form
\[
  m_{\rm arith}(z)
  =A+\int_{\mathbb R}
  \left(\frac1{t-z}-\frac{t}{1+t^2}\right)\,d\mu_{\rm arith}(t),
  \qquad
  A=\Re m_{\rm arith}(i)\in\mathbb R,
\]
geschrieben. Die äquivalente Nullstellensumme ist im kanonisch symmetrischen
Sinn zu verstehen.
\end{definition}

\begin{proposition}[Verschwinden des linearen Nevanlinna-Terms]
\label{prop:b-zero}
Die allgemeine Herglotz--Nevanlinna-Darstellung erlaubt einen zusätzlichen
Term $bz$ mit $b\ge0$. Für das hier betrachtete spezielle
$m_{\rm arith}=-\Xi'/\Xi$ gilt jedoch $b=0$. Wegen der symmetrischen
$\pm\gamma$-Paarung und der Hadamard-Darstellung konvergiert
\[
  \sum_{\gamma>0}m_\gamma
  \left(\frac1{\gamma-z}+\frac1{-\gamma-z}\right)
\]
unter Verwendung von
\[
  \sum_\gamma\frac{m_\gamma}{\gamma^2}<\infty
\]
ohne zusätzlichen linearen Anteil.
\end{proposition}

\begin{theorem}[Herglotz-Kriterium]\label{thm:herglotz-rh}
Für die arithmetische Funktion gilt
\[
  m_{\rm arith}\text{ ist Herglotz}
  \quad\Longleftrightarrow\quad
  \RH.
\]
\end{theorem}

\begin{definition}[Normierte Weil-Distribution]\label{def:weil-distribution}
Nullstellenseite und arithmetische Seite sind zwei Darstellungen derselben
normierten Distribution:
\[
  W_\xi^{\rm norm}
  =W_{\rm zeros}
  =W_{\rm pole/triv}+W_\Gamma+W_{\rm prime}.
\]
Sie sind nicht miteinander zu addieren.
\end{definition}

\begin{theorem}[Kanonische Positivierungsbrücke]\label{thm:positivity-bridge}
Mit Involution und Faltungskonvention aus P02 liefert
\[
  \Phi=\phi^\ast *\phi
\]
die Identität
\[
  W_\xi^{\rm norm}[\Phi]
  =Q_{\rm zeros}[\phi]
  =\sum_{\gamma\in\Gamma}
    m_\gamma\,|\widehat\phi(\gamma)|^2.
\]
Dies ist die im Programm verwendete kanonische Positivierungsbrücke.
Ihre Einzigkeit ist nicht bewiesen.
\end{theorem}

\begin{theorem}[Interface-Schutzsatz]\label{thm:interface-firewall}
Die lineare und die quadratische Ebene bleiben typgetrennt:
\[
\begin{aligned}
  m_{\rm arith}&\rightsquigarrow W_\xi^{\rm norm},\\
  a,b&\longmapsto g_{a,b}\longmapsto h_{a,b}
       \longmapsto B_W(a,b).
\end{aligned}
\]
Ohne Autokorrelations-/Positivierungspaarung wird
$m_{\rm arith}$ nicht direkt mit der quadratischen Weilform identifiziert.
\end{theorem}

\begin{openproblem}[Rückbindung an Objekt X]\label{op:pi-gamma}
Die vorgeschlagene Identifikation
\[
  m_{\rm arith}=\Pi_\gamma(X)
\]
bleibt offen.
\end{openproblem}

% ----------------------------------------------------------------
\section{Jacobi-/Herglotz-Realisierung: konditionale Architektur}
% ----------------------------------------------------------------

\begin{definition}[Spektralmaß bei selbstadjungierter Realisierung]
\label{def:spectral-measure}
Angenommen, $A_N^{\rm Jac,-}$ besitzt eine selbstadjungierte Realisierung, und
$\Omega_N$ sei normiert, $\|\Omega_N\|=1$. Dann definiert man
\[
  \mu_{\Omega,N}(B)
  :=\langle\Omega_N,E_{A_N}(B)\Omega_N\rangle
\]
und
\[
  m_{\Omega,N}(z)
  :=\langle\Omega_N,(A_N^{\rm Jac,-}-z)^{-1}\Omega_N\rangle
  =\int_{\mathbb R}\frac{d\mu_{\Omega,N}(\lambda)}{\lambda-z}.
\]
Damit gilt
\[
  \mu_{\Omega,N}(\mathbb R)=1.
\]
\end{definition}

\begin{openproblem}[Selbstadjungiertheit des Jacobi-Kandidaten]
\label{op:jacobi-selfadjoint}
Für den konkreten Kandidaten
\[
  A_N^{\rm Jac,-}=H_N+\beta_NJ_N^-
\]
ist eine selbstadjungierte Realisierung nicht bewiesen.
\end{openproblem}

\begin{definition}[Nevanlinna-normalisierte Approximanten]
\label{def:renorm-approximants}
Da $\mu_{\Omega,N}(\mathbb R)=1$, während
$\mu_{\rm arith}$ unendliche Gesamtmasse besitzt, ist für einen möglichen
Grenzübergang eine positive Skalierung einzuführen:
\[
  \widetilde\mu_N:=c_N\mu_{\Omega,N},
  \qquad c_N>0.
\]
Die Nevanlinna-kompatible Form lautet
\[
  \widetilde m_N^{\rm ren}(z)
  :=a_N+\int_{\mathbb R}
  \left(\frac1{t-z}-\frac{t}{1+t^2}\right)\,d\widetilde\mu_N(t),
  \qquad a_N\in\mathbb R.
\]
Nur unter zusätzlicher Symmetrie reduziert sich dies auf eine reine Skalierung
$c_Nm_{\Omega,N}$.
\end{definition}

\begin{theorem}[Konditionale Herglotz-Firewall]
\label{thm:conditional-firewall}
Falls echte Herglotzfunktionen lokal gleichmäßig in $\mathbb C^+$ gegen
$m_{\rm arith}$ konvergieren,
\[
  \widetilde m_N^{\rm ren}
  \xrightarrow[N\to\infty]{\rm loc.\,unif.}
  m_{\rm arith},
\]
und die nichtdegenerierte Herglotzstruktur im Grenzübergang erhalten bleibt,
dann ist $m_{\rm arith}$ Herglotz und damit gilt RH:
\[
  \widetilde m_N^{\rm ren}\to m_{\rm arith}
  \quad\Longrightarrow\quad
  m_{\rm arith}\text{ Herglotz}
  \quad\Longrightarrow\quad
  \RH.
\]
\end{theorem}

\begin{openproblem}[Daten für den Jacobi-/Herglotz-Grenzübergang]
\label{op:jacobi-limit}
Der Grenzübergang selbst bleibt offen. Insbesondere fehlen:
\begin{enumerate}
  \item eine korrekte selbstadjungierte Realisierung von
        $A_N^{\rm Jac,-}$;
  \item eine kanonische Wahl von $\Omega_N$;
  \item eine kanonische Skalierung $c_N>0$ und Kontrolle der Konstanten
        $a_N\in\mathbb R$;
  \item vague/lokale Maßkonvergenz
        $\widetilde\mu_N\to\mu_{\rm arith}$;
  \item Kontrolle der Nevanlinna-Gewichte, etwa von
        \[
          \int_{\mathbb R}
          \frac{d\widetilde\mu_N(t)}{1+t^2};
        \]
  \item Tail-Kontrolle, die aus der Maßkonvergenz die benötigte lokal
        gleichmäßige Konvergenz der Herglotzfunktionen ableiten lässt.
\end{enumerate}
Vague Konvergenz allein wird ausdrücklich nicht als hinreichend behauptet.
\end{openproblem}

\begin{remark}[Konditionale Architektur]\label{rem:conditional-chain}
Die offene Architektur lautet
\[
  A_N^{\rm Jac,-}
  \stackrel{?[O]}{\longrightarrow}
  m_{\Omega,N}
  \stackrel{c_N,a_N}{\longrightarrow}
  \widetilde m_N^{\rm ren}
  \stackrel{?[O]}{\longrightarrow}
  m_{\rm arith}.
\]
\end{remark}

% ----------------------------------------------------------------
\section{Statusmatrix}
% ----------------------------------------------------------------

\small
\begin{longtable}{>{\raggedright\arraybackslash}p{0.55\textwidth}
                  >{\raggedright\arraybackslash}p{0.21\textwidth}
                  >{\raggedright\arraybackslash}p{0.15\textwidth}}
\toprule
Aussage & Status & Quelle\\
\midrule
\endfirsthead
\toprule
Aussage & Status & Quelle\\
\midrule
\endhead
$D_N(z)\to e^{-\gamma^2/4}$ (Determinanten-No-Go)
& PROVED & NEU-091\\
Quadratischer Pivot $Q_N$
& PROVED (Definition) & NEU-091\\
Testkegel-Positivität $Q_N[\phi]\ge0$ auf $\mathcal C$
& PROVED & NEU-092\\
Bilinearer Lift $B_N(f,g)$ positiv-semidefinit auf $\mathcal S$
& OPEN & NEU-092\\
Abstrakte PSD-Architektur $K_N$ aus $\rho_N$
& PROVED & NEU-093\\
Logarithmischer Bochner-Kandidat $\Psi_N$
& PROVED & NEU-094\\
Kanonische Wahl von $\rho_N$
& OPEN & NEU-093\\
Skalentriage: $T\gg M_N$ Diagonal-Kollaps
& PROVED & NEU-096\\
Skalentriage: $T=O(1)$ Rang-eins-Kollaps
& PROVED & NEU-096\\
Zwischenregime $1\ll T\ll M_N$ Weil-tauglich
& OPEN & NEU-096/097\\
Hardy--Littlewood Singulärserien-Hauptterm
& CONDITIONAL & NEU-098\\
$\mathcal V(M,H)\sim H\log(M/H)$ uniform in
$1\le H\le M^{1-\varepsilon}$
& CONDITIONAL (RH+SPC) & NEU-101\\
$\mathcal V(M,\sqrt M)\sim\frac12\sqrt M\log M$
& CONDITIONAL (RH+SPC) & NEU-101\\
$\mathcal P^{\rm unf}_{N,H}$ lokales Formfaktor-Ersatzobjekt
& PROVED$_{\rm part}$ & NEU-104\\
$\LFF\Rightarrow$ Rampenasymptotik
& PROVED$_{\rm part}$ & NEU-107\\
$\LFF\Leftrightarrow$ Rampenform
& NO-GO; nur $\Rightarrow$ & NEU-107\\
LFF allein identifiziert $Q_{\rm Weil}$ nicht
& PROVED$_{\rm part}$ (Typwarnung) & NEU-108\\
$\sigma_{\rm loc}(Q_{\rm Weil})\stackrel?=|\alpha|$
& OPEN & NEU-110\\
$m_{\rm arith}$ Herglotz $\Leftrightarrow$ RH
& PROVED & NEU-111\\
$\mu_{\rm arith}=\sum m_\gamma\delta_\gamma$ (reines Nullstellenmaß)
& PROVED & NEU-112/118\\
$m_{\rm arith}$ in Nevanlinna-Form (symmetrische Konvergenz)
& PROVED (Typisierung) & NEU-111/118\\
Linearer Nevanlinna-Koeffizient $b=0$
& PROVED & NEU-111/118\\
$W_\xi^{\rm norm}=W_{\rm zeros}$ (keine Doppelzählung)
& PROVED & NEU-113/115\\
Kanonische Positivierungsbrücke
$Q_{\rm zeros}=\sum m_\gamma|\widehat\phi|^2$
& PROVED & NEU-112/113\\
Einzigkeit der Positivierungsbrücke
& OPEN & ---\\
$m_{\rm arith}=\Pi_\gamma(X)$
& OPEN & NEU-114\\
$A_N^{\rm Jac,-}$ selbstadjungiert
& OPEN & NEU-119\\
Nevanlinna-Renormierung $(c_N,a_N)$ und Tail-Kontrolle
& OPEN & NEU-120\\
Vague Konvergenz $\widetilde\mu_N\to\mu_{\rm arith}$
& OPEN & NEU-120\\
$\widetilde m_N^{\rm ren}\to m_{\rm arith}\Rightarrow\RH$
& CONDITIONAL & NEU-120\\
$\mu_{\Omega,N}(\mathbb R)=1$ ohne $c_N$ genügt
& NO-GO & NEU-120\\
Pole $\pm i/2$ in $m_{\rm arith}$
& NO-GO & NEU-120\\
\bottomrule
\end{longtable}
\normalsize

% ----------------------------------------------------------------
\section{Provenienz und Manuskriptstatus}
% ----------------------------------------------------------------

\begin{remark}[Kompakte Provenienz]\label{rem:provenance}
Dieses Manuskript konsolidiert NEU-091--120 unter
\texttt{PASS-A-PROTOKOLL.md}. Direkte Markdown-Quelle ist
\texttt{papers/P07\_Weil\_Form\_Statistics.md} im Status
\emph{SYN FINAL AUDITED}, Commit \texttt{6a162f92}; die
Goldston--Montgomery-Synchronisierung von NEU-101 liegt in Commit
\texttt{92d731d1}. P02 liefert die hier referenzierten kanonischen
Fourier-, Paley--Wiener- und Weilform-Konventionen. P03 wurde nur für
SYN-Stil und Konventionskonsistenz herangezogen; es werden daraus keine
zusätzlichen P07-Aussagen importiert.
\end{remark}

\begin{remark}[LaTeX-SYN-Finalaudit]\label{rem:audit-status}
Der LaTeX-Transferaudit ist abgeschlossen. Geprüft wurden ausschließlich
Markdown--LaTeX-Übereinstimmung, Formeltranskription,
Fourier-/Mellin-Konventionen, Satzstatus,
Definition-versus-Hypothese-versus-Open-Problem, Referenzen, Labels,
Doppelzählungen, unzulässige Hochstufungen und Kompilierbarkeit.
Es wurde kein neuer mathematischer Konflikt gefunden und kein Forschungsknoten
NEU-091--120 wieder geöffnet. Die in der Markdown-Fassung typographisch
verkürzte Interface-Kette wurde gemäß der verbindlichen P02-Typisierung als
$(a,b)\mapsto g_{a,b}\mapsto h_{a,b}\mapsto B_W(a,b)$ geschrieben.
Damit trägt diese Fassung den internen Status
\[
  \boxed{\texttt{P07\_Weil\_Form\_Statistics.tex}\text{ --- SYN FROZEN }\checkmark[K/M]}.
\]
\end{remark}

\begin{thebibliography}{9}

\bibitem{Bombieri2000}
E.~Bombieri,
\emph{Remarks on Weil's quadratic functional in the theory of prime numbers},
2000.

\bibitem{Connes1999}
A.~Connes,
\emph{Trace formula in noncommutative geometry and the zeros of the Riemann
zeta function},
1999.

\bibitem{GoldstonMontgomery1987}
D.~A.~Goldston and H.~L.~Montgomery,
\emph{Pair correlation of zeros and primes in short intervals},
in \emph{Analytic Number Theory}, 1987, pp.~183--203.

\bibitem{Akhiezer1965}
N.~I.~Akhiezer,
\emph{The Classical Moment Problem},
1965.

\bibitem{Simon2011}
B.~Simon,
\emph{Szeg\H{o}'s Theorem and Its Descendants},
2011.

\bibitem{HardyLittlewood1923}
G.~H.~Hardy and J.~E.~Littlewood,
\emph{Some problems of Partitio Numerorum III},
1923.

\end{thebibliography}

\end{document}
